% Generated draft: requires mathematical and source review.
\documentclass[11pt]{article}
\usepackage[margin=25mm]{geometry}
\usepackage{fontspec}
\setmainfont{Latin Modern Roman}
\usepackage{amsmath,amssymb,mathtools,unicode-math}
\setmathfont{Latin Modern Math}
\usepackage{enumitem,needspace,xurl,hyperref}
\setlength{\parindent}{0pt}
\setlength{\parskip}{5pt}
\setlength{\emergencystretch}{3em}
\newcommand{\N}{\mathbb{N}}
\newcommand{\Z}{\mathbb{Z}}
\newcommand{\Q}{\mathbb{Q}}
\newcommand{\R}{\mathbb{R}}
\newcommand{\C}{\mathbb{C}}
\newcommand{\F}{\mathbb{F}}
\newcommand{\A}{\mathbb{A}}
\newcommand{\PP}{\mathbb{P}}
\newcommand{\E}{\mathbb{E}}
\newcommand{\eps}{\varepsilon}
\newcommand{\dd}{\,\mathrm{d}}
\begin{document}
\section*{UnsolvedMath ID: 30003348}
\textbf{Source problem code:} catalogue entry
\subsection{Definitions and mathematical statement}
Let \(n>2\). For a smooth Riemannian \(n\)-manifold \((M,g)\), write \(SM=\{v\in TM:g(v,v)=1\}\) and let \(\varphi_g^t\) denote its unit-speed geodesic flow, assumed defined for every \(t\in\R\). The flow is Anosov if there is a continuous invariant splitting
\[
T(SM)=E^s\oplus\R X_g\oplus E^u
\]
with constants \(C\geq1\) and \(\lambda>0\) such that, for all \(t\geq0\),
\[
\|D\varphi_g^t(v)\|\leq Ce^{-\lambda t}\|v\|\quad(v\in E^s),\qquad
\|D\varphi_g^{-t}(v)\|\leq Ce^{-\lambda t}\|v\|\quad(v\in E^u),
\]
where \(X_g\) is the generating vector field and the norms come from a fixed auxiliary Riemannian metric on \(SM\). The source does not specify whether \(M\) must be closed, connected, or subject to additional geometric hypotheses; these are unresolved parts of the intended domain.

A periodic orbit of period \(T>0\) determines a closed geodesic of length \(T\). A marking identifies the free homotopy classes used to compare closed geodesics on the manifolds under consideration. When each relevant nontrivial free homotopy class \(c\) has a unique closed geodesic representative \(\gamma_c\), the marked length spectrum is the function
\[
L_g(c)=\operatorname{length}_g(\gamma_c).
\]
The domain, marking, and any hypotheses ensuring this description are not supplied. If uniqueness is unavailable, a length set or multiset for each class would require a further convention. The phrase ``periodic data moduli'' in the original statement, rewritten as ``periodic-orbit data'' in the supplied curated statement, is not defined by the record.

\textbf{Problem as supplied.} In dimensions greater than two, are periodic-orbit data and the marked length spectrum complete invariants of Anosov geodesic flows?

To specify completeness precisely, one must fix a class \(\mathcal A_n\) of admissible marked geometric objects, an equivalence relation \(\sim\), a periodic-data invariant \(P\), and the marked length invariant \(L\). Joint completeness would then mean
\[
\forall a,b\in\mathcal A_n,\qquad
\bigl(P(a)=P(b)\ \text{and}\ L(a)=L(b)\bigr)\Longrightarrow a\sim b.
\]
If the question instead asks whether each invariant is complete separately, the required implications are
\[
\forall a,b\in\mathcal A_n,\qquad
P(a)=P(b)\Longrightarrow a\sim b,
\qquad
L(a)=L(b)\Longrightarrow a\sim b.
\]
These are distinct interpretations. The supplied record does not select between them, define \(P\), or specify whether \(\sim\) means metric isometry, flow conjugacy with a stated regularity, orbit equivalence, or another relation. Thus the displayed implications explain the missing specifications rather than assert a uniquely determined conjecture.
\subsection{Short English statement}
The supplied question asks whether periodic-orbit information and the lengths of closed geodesics indexed by a marking completely determine an Anosov geodesic flow in dimension \(n>2\). A precise conjecture still requires the admissible manifolds, the meaning of periodic data, the marking, the intended equivalence, and whether completeness is joint or separate. The supplied open status is dated curation, not an independently verified conclusion.
\subsection{Sources}
\begin{itemize}
\item Local UnsolvedMath registry, permanent ID \(30003348\), title ``Completeness of Periodic Data and Marked Length Spectra,'' source problem code \texttt{catalogue entry}, in the Oberwolfach Reports Open Problems collection. The original wording is ``In higher dimensions, are periodic data moduli and the marked length spectrum complete invariants?'' No individual proposer is supplied; the proposed year is recorded as 2017.
\item Supplied source citation: \emph{Mini-Workshop: Spaces and Moduli Spaces of Riemannian Metrics (2018)}, Oberwolfach Reports, pp.\ 133--166, DOI \url{https://doi.org/10.4171/owr/2017/3}. This link is also listed among the literature sources as ``Oberwolfach report on rigidity questions.'' The supplied citation's 2018 designation and its 2017 DOI component are retained without resolving their relationship.
\item Supplied literature source: ``The marked length spectrum of Anosov manifolds,'' \url{https://doi.org/10.4007/annals.2019.190.1.6}.
\item Supplied literature source: ``Marked length pattern rigidity for arithmetic manifolds,'' \url{https://doi.org/10.1017/etds.2024.72}.
\item Literature assessment checked at \texttt{2026-08-21}: the supplied curation labels the problem open and reports strong local and special-case rigidity results for periodic data and marked length spectra, including all-dimensional local marked-length rigidity under nonpositive curvature, while describing completeness in the unrestricted higher-dimensional setting as unresolved. This is supplied dated curation derived from the OWR record, not independently verified literature assessment; the status and cited claims require independent verification.
\item The registry labels the statement \texttt{corrected\_verified} and reports that the official passage was checked and rewritten while preserving its mathematical content. This verification claim belongs to the supplied record; no source was independently consulted for this draft.
\end{itemize}
\end{document}
